\documentclass[11pt]{article}
\usepackage[margin=1in]{geometry}
\usepackage{amsmath,amssymb,amsthm,graphicx,microtype}
\usepackage[hidelinks]{hyperref}
\newtheorem{theorem}{Theorem}
\newtheorem{proposition}{Proposition}
\newtheorem{lemma}{Lemma}
\newcommand{\A}{\mathcal A}
\newcommand{\M}{M}
\newcommand{\C}{\mathbb C}
\newcommand{\od}{\mathbin{\odot}}
\newcommand{\tmax}{\otimes_{\max}}
\newcommand{\tmin}{\otimes_{\min}}
\setlength{\parindent}{0pt}
\setlength{\parskip}{0.65em}

\title{Maximal-Tensor Multiplication Does Not Extend\\
on the CAR Algebra}
\author{Run \texttt{fa\_banach\_001}, lane 8}
\date{22 July 2026}

\begin{document}
\maketitle

\textbf{Status.} Candidate counterexample, likely valid, for human review.
The proof is complete; novelty is supported only by the bounded search in
Section~\ref{sec:novelty}.

\section{Source question}

Let $A$ be a unital $C^*$-algebra.  Lesch writes
\[
 \mu_n:A^{\od(n+1)}\longrightarrow A,
 \qquad
 \mu_n(a_0\otimes\cdots\otimes a_n)=a_0a_1\cdots a_n,
\]
and asks whether it extends continuously to the maximal $C^*$-tensor product
$A^{\tmax(n+1)}$ \cite[Section~3.1]{Lesch}.  He says that he suspects
interesting examples where it does not extend, but that an expert poll was
inconclusive.  Liu later describes the extension question as ``still
inconclusive'' and cites Lesch \cite[Section~2, PDF p.~5]{Liu}.

\begin{center}
\includegraphics[width=\textwidth]{figures/open_problem_crop.png}
\end{center}

Liu's displayed multiplication map uses a shifted number of tensor factors.
To avoid this notational issue, below $m_k$ always means multiplication on
exactly $k$ algebraic factors.  Thus Lesch's $\mu_n$ is $m_{n+1}$.

The same sources consider the contraction map
\[
 \iota:A^{\od(n+1)}\longrightarrow
 L\bigl(A^{\od n},A\bigr),
\]
where
\[
 \iota(a_0\otimes\cdots\otimes a_n)
 (b_1\otimes\cdots\otimes b_n)
 =a_0b_1a_1\cdots b_na_n.
\]
They ask whether a suitable joint maximal-tensor extension exists.

\section{Definitions}

For $C^*$-algebras $B$ and $D$, $B\od D$ denotes their algebraic tensor
product, while $B\tmin D$ and $B\tmax D$ denote the minimal and maximal
$C^*$-completions.  The minimal tensor product is injective: inclusions of
$C^*$-algebras induce isometric inclusions of minimal tensor products.  A
$C^*$-algebra $A$ is nuclear if the canonical maximal-to-minimal quotient is
an isomorphism for $A\od B$ and every $C^*$-algebra $B$; see, for example,
\cite[Chapters 2--3]{BrownOzawa}.

The CAR algebra is
\[
 \A=\M_{2^\infty}
   =\varinjlim_r \M_2^{\otimes r}.
\]
It is unital, separable, simple, and nuclear.  Its $r$th finite stage is a
unital copy of $\M_d$, where $d=2^r$.

\section{Proof intuition}

Matrix multiplication hides a dimension factor that the spatial tensor norm
does not see.  In $\M_d\od\M_d$, the tensor
$z_d=\sum_j e_{1j}\otimes e_{j1}$ merely moves one $d$-dimensional coordinate
subspace isometrically to another, so its tensor norm is one.  Multiplication,
however, collapses all $d$ summands onto the same matrix unit and produces
$d e_{11}$.  The CAR algebra contains such matrix blocks for unbounded $d$.
Nuclearity is the bridge from this spatial calculation to the maximal tensor
norm asked about in the source.

\section{Counterexample}

We first record the reusable obstruction.

\begin{proposition}[Unbounded matrix-block criterion]\label{prop:criterion}
Let $A$ be a unital nuclear $C^*$-algebra.  Suppose that $A$ contains
$C^*$-subalgebras isomorphic to $\M_d$ for an unbounded set of integers $d$.
Then, for every $k\geq2$, the algebraic multiplication map
\[
 m_k:A^{\od k}\longrightarrow A,
 \qquad
 m_k(a_1\otimes\cdots\otimes a_k)=a_1\cdots a_k,
\]
is unbounded for the maximal $C^*$-tensor norm.  Hence it has no continuous
extension to $A^{\tmax k}$.
\end{proposition}

\begin{proof}
Fix one of the subalgebras $\M_d\subseteq A$ and choose standard matrix units
$(e_{ij})_{i,j=1}^d$.  In $\M_d\od\M_d$ define
\[
 z_d=\sum_{j=1}^d e_{1j}\otimes e_{j1}.
\]
Let $(\xi_j)_{j=1}^d$ be the standard basis of $\C^d$.  In the spatial
representation on $\C^d\otimes\C^d$,
\[
 z_d(\xi_p\otimes\xi_q)
 =\begin{cases}
    \xi_1\otimes\xi_p,&q=1,\\
    0,&q\ne1.
  \end{cases}
\]
Thus $z_d$ is a partial isometry from
$\C^d\otimes\C\xi_1$ onto $\C\xi_1\otimes\C^d$, and
\begin{equation}\label{eq:zdnorm}
 \|z_d\|_{\M_d\tmin\M_d}=1.
\end{equation}

By injectivity of the minimal tensor product, the inclusion
$\M_d\subseteq A$ preserves the norm in \eqref{eq:zdnorm}.  Since $A$ is
nuclear, maximal and minimal tensor powers of $A$ coincide.  Consequently,
\[
 \|z_d\|_{A\tmax A}=1.
\]
But binary multiplication gives
\begin{equation}\label{eq:blowup}
 m_2(z_d)=\sum_{j=1}^d e_{1j}e_{j1}=d e_{11},
 \qquad \|m_2(z_d)\|=d.
\end{equation}
As the available dimensions are unbounded, $m_2$ is unbounded.

For $k>2$, put
\[
 z_d^{(k)}=z_d\otimes 1_A^{\otimes(k-2)}\in A^{\od k}.
\]
Tensoring by identities preserves the $C^*$-norm, so
$\|z_d^{(k)}\|_{A^{\tmax k}}=1$, whereas
$m_k(z_d^{(k)})=d e_{11}$.  Hence every $m_k$, $k\ge2$, is unbounded.
\end{proof}

\begin{theorem}[CAR counterexample]\label{thm:main}
For the CAR algebra $\A=\M_{2^\infty}$ and every $n\ge1$, Lesch's map
\[
 \mu_n:\A^{\od(n+1)}\longrightarrow\A
\]
does not extend continuously to $\A^{\tmax(n+1)}$.  Consequently, the
contraction map in Liu and Lesch has no continuous extension
\[
 \iota:\A^{\tmax(n+1)}\longrightarrow
 L_{\mathrm{cont}}(\A^{\tmax n},\A)
\]
that agrees with the algebraic contraction.
\end{theorem}

\begin{proof}
The CAR algebra is nuclear and contains the unital finite stages
$\M_{2^r}$ for all $r$.  Proposition~\ref{prop:criterion}, with
$k=n+1$, proves the first assertion.

For the second, suppose such an extension of $\iota$ existed.  Evaluation at
$1_A^{\otimes n}$ is a bounded map from
$L_{\mathrm{cont}}(\A^{\tmax n},\A)$ to $\A$.  On algebraic tensors its
composition with $\iota$ is exactly
\[
 a_0\otimes\cdots\otimes a_n\longmapsto a_0a_1\cdots a_n=\mu_n(a).
\]
The composition would therefore give the forbidden bounded extension of
$\mu_n$.
\end{proof}

This supplies a separable, unital, simple, nuclear example.  It also shows
that the failure is simultaneous in every nontrivial multilinear degree.

\section{Verification and limitations}

The script \texttt{code/verifier.py} constructs $z_d$ as an operator on
$\C^d\otimes\C^d$ and checks, for $d=2,4,8,16$, that
\[
 \|z_d\|=1,\qquad z_dz_d^*z_d=z_d,
 \qquad m_2(z_d)=d e_{11}.
\]
It runs with:
\begin{verbatim}
conda run --no-capture-output -n sandbox python code/verifier.py
\end{verbatim}
All residuals are zero and the computed multiplication norms are
$2,4,8,16$.  This check guards against an index or tensor-order error; the
proof for all $d=2^r$ is the displayed matrix-unit calculation.

The theorem answers the extension question negatively in general.  It does
not classify all $C^*$-algebras for which multiplication is bounded, and it
does not settle the case of the irrational rotation algebra occurring in
Liu's geometric applications.

\section{Bounded novelty check}\label{sec:novelty}

The search on 22 July 2026 covered the run's registry, solution, attempt, and
proof-gap indexes; the exact source ids and titles; and web/arXiv searches
combining ``multiplication map'' and ``maximal $C^*$-tensor product'' with
CAR, UHF, nuclear, minimal tensor product, and subhomogeneous.

Vaes and Van Daele state that multiplication generally need not extend to the
minimal $C^*$-tensor product \cite[Introduction]{VaesVanDaele}.  Van Daele's
later survey points to growth of the multiplication norms of $\M_d$ as
$d\to\infty$ \cite[Remark~3.3]{VanDaele}.  These are close precursors, but the
inspected passages do not state the fixed CAR/UHF example or transfer it to
the maximal norm through nuclearity.  A 2021 MathOverflow discussion gives a
spatial-tensor obstruction for $B(H)$, which is not by itself a maximal-norm
answer.  No explicit CAR/UHF maximal-tensor counterexample was found.

This supports, but does not certify, novelty.  Because the construction is an
elementary synthesis of standard facts, the priority claim requires a deeper
operator-algebra literature review.  The mathematical validity of the
counterexample does not depend on that novelty assessment.

\section{Human-review recommendation}

Send the packet to an operator-algebra or $C^*$-tensor-product specialist.
The proof audit should focus on the injective minimal-tensor inclusion and the
nuclear maximal-to-minimal identification.  The larger review burden is a
specialist novelty search, particularly in the Hopf-$C^*$-algebra and
operator-space literature.

\begin{thebibliography}{9}
\bibitem{Liu}
Y. Liu,
\emph{Hypergeometric function and Modular Curvature I. Hypergeometric
functions in Heat Coefficients}, arXiv:1810.09939 (2018),
\url{https://arxiv.org/abs/1810.09939}.

\bibitem{Lesch}
M. Lesch,
\emph{Divided differences in noncommutative geometry: rearrangement lemma,
functional calculus and expansional formula},
J. Noncommut. Geom. 11 (2017), 193--223; arXiv:1405.0863,
\url{https://arxiv.org/abs/1405.0863}.

\bibitem{BrownOzawa}
N. P. Brown and N. Ozawa,
\emph{$C^*$-Algebras and Finite-Dimensional Approximations},
Graduate Studies in Mathematics 88, American Mathematical Society, 2008.

\bibitem{VaesVanDaele}
S. Vaes and A. Van Daele,
\emph{Hopf $C^*$-algebras},
Proc. London Math. Soc. 82 (2001), 337--384; arXiv:math/9907030,
\url{https://arxiv.org/abs/math/9907030}.

\bibitem{VanDaele}
A. Van Daele,
\emph{From Hopf algebras to topological quantum groups: a short history,
various aspects and some problems},
Banach Center Publ. 120 (2020), 179--223,
doi:10.4064/bc120-14.
\end{thebibliography}

\end{document}
